% Provide only the theory needed to understand Chapter~\ref{chap:methodology}.
% Reuse a small running example where it improves clarity.

\section{Routing problems}\label{sec:routing-background}

\subsection{Traveling salesman problem}
% Define graph, nodes, arcs, tour, and travel cost briefly.

\subsection{Vehicle routing problem}
% Explain the transition from one tour to a fleet of capacitated routes.

\subsection{Relevant VRP variants}
% Introduce multiple trips, multiple products, rich-routing features, and
% sequence-dependent costs. State explicitly what ``trip'' means in this work.

\section{Optimization methods}\label{sec:optimization-background}

\subsection{Mathematical programming}
% Introduce variables, objective functions, constraints, and MILP. Distinguish
% an exact formulation from a solver run that actually proves optimality.

\subsection{Heuristics}
% Explain construction and improvement heuristics, then best-fit packing,
% nearest neighbor, and open-path 2-opt as they are used later.

\section{Constraint programming}\label{sec:cp-background}

\subsection{Basic concepts}
% Variables, domains, constraints, propagation, search, and optimization. Add a
% small domain-reduction example.

\subsection{Routing models}
% Briefly introduce successor/circuit and other common representations.

\subsection{Sequence variables}
% Explain membership, ordering, possible insertions, and endpoints with one
% partial insertion-based sequence. Do not confuse this representation with
% similarly named scheduling constructs.
